% ============================================================================
% PART II, CH B — INTERPRETABILITY KNOBS AND METRICS      [slug: knobs-metrics]
% Second of the three Part II chapters (split 2026-07-14): the VOCABULARY.
% No routes named yet (they are chapter C's); the differentiability gate
% appears only as a per-row property ("has a plausible surrogate or not").
% ALL experiments live in chapter C (Matteo, 2026-07-14) — this chapter is
% taxonomy only.
% STATUS: taxonomy artifact exists (Master/docs/
% interpretability_knobs_metrics_taxonomy.md); open thread: merge Matteo's
% own knob/metric list into it.
% Material map: ../outline.md#knobs-metrics
% ============================================================================
\cleardoublepage
\chapter{Interpretability Knobs and Metrics (5--7)}   % page goal — scaffolding, DELETE before submission
\label{ch:iq-knobs}

\section{Knobs and Metrics: Two Sides of the Same Property}
\label{sec:iq-knob-metric-axis}
\vault{2026-07-12_1410_knob-metric-axis-and-principledness-spectrum}
% The founding distinction (vault 2026-07-12_1410, idea 1): a knob is an
% INPUT to training (cause side), a metric is an OUTPUT measured on a
% trained model (effect side) — most interpretability properties can appear
% as both, and "diversity" is the classic trap (penalty weight AND measured
% spread). The bridge: a knob is usually the differentiable training-time
% SURROGATE of a metric you care about at evaluation time. Consequence
% stated here without route names: properties split into "has a plausible
% differentiable surrogate" vs "doesn't" — and empirically, sorting our own
% list, purity / cosine-readout reliability (Steck) / seed-stability have
% blank surrogate columns (structurally measure-and-select only; chapter C
% picks this up as the route decider).

\section{A Taxonomy of Knobs and Metrics}
\label{sec:iq-taxonomy}
\vault{2026-07-12_1411_interpretability-knobs-map-to-rudin-families, 2026-07-12_1413_understandability-weighted-feature-use-and-sparsity-axes, 2026-07-13_2326_part-ii-framework-knob-entry-points-explanation-backward-tracing, 2026-07-12_1410_knob-metric-axis-and-principledness-spectrum}
% The property -> metric -> surrogate/knob table (artifact: Master/docs/
% interpretability_knobs_metrics_taxonomy.md). Content blocks:
%  1. Rudin's constraint families as the citable spine (vault
%     2026-07-12_1411): her master enumeration is explicitly open-ended and
%     domain-specific — the licence to add recsys-specific rows under her
%     families. Free-citation mappings: separation ~ disentanglement
%     (R22 §6 verbatim), popularity-decorrelation ~ concept whitening,
%     non-negativity ~ half of monotonicity. Plus the soft-penalty vs
%     hard-constraint tag (R22 Eq.1) on every row.
%  2. Entry-point classification (Matteo's taxonomy, vault 2026-07-13_2326):
%     (a) outside the model (hyperopt-only, accuracy window — e.g.
%     n_prototypes; open empirical sub-question: is "more prototypes" a
%     consistent positive direction? test, don't assume), (b) measurable
%     properties affected but not trained for (sharpness/distinctness —
%     the blind-tuning mismatch as a taxonomy class), (c) loss-enterable
%     (the differentiability gate).
%  3. Sparsity is FOUR distinct axes (vault 2026-07-12_1413): activation
%     sparsity / prototype-feature sparsity / model-size sparsity /
%     feature-set sparsity — each with its own metric and knob; crispness is
%     NOT prototype-feature sparsity (crisp-but-dense is possible). Carry
%     Rudin's caveat: sparsity != interpretability.
%  4. The understandability-weighted feature-use row is IN the taxonomy,
%     with its NOVELTY highlighted (Matteo, 2026-07-14): it instantiates the
%     VOCABULARY dimension (bullets below); the full mechanism treatment
%     stays in the tuning chapter.
%  5. OPEN THREAD: merge Matteo's own knob/metric list into the taxonomy.
% THE VOCABULARY DIMENSION (Matteo, 2026-07-14 — a novelty flag of the
% taxonomy). C6 citation clarified (Matteo): cite the ORIGINAL Niknazar &
% Mednick academic publication — the poster in the literature folder is
% Matteo's own poster made FOR that publication, not the citable object; pull
% the paper's full bib data at writing time.
\begin{itemize}
    \item In a truly interpretable model, explanation vocabulary $\approx$ internal vocabulary. \vault{2026-07-14_1154_vocabulary-dimension-constraining-vocab-interpretability}
    \item Constraining a model's vocabulary = a rarely-discussed interpretability dimension (citation: Niknazar \& Mednick, C6). \vault{2026-07-14_1154_vocabulary-dimension-constraining-vocab-interpretability}
    \item Vocabulary constraint doubles as expert-knowledge infusion — and can even increase performance. \vault{2026-06-11_1555_interpretability-constraints-as-expert-knowledge-infusion}
    \item The UW feature-use row instantiates the vocabulary dimension — novel (Rudin's principle, our mechanism). \vault{2026-07-12_1413_understandability-weighted-feature-use-and-sparsity-axes}
\end{itemize}
% ROLE FOR THE NEXT CHAPTER (Matteo, 2026-07-14): the taxonomy is the LOOKUP
% TABLE the tuning method consults — per knob it must answer (a) which
% metric(s) quantify the property, (b) the ADJUSTMENT CLASS (bullets below,
% higher abstraction level, details to be determined). The method (next
% chapter) is the procedure that dispatches on exactly this classification.
\begin{itemize}
    \item Per knob, the taxonomy carries its \emph{adjustment class}: model-structural vs training-adjustable. \vault{2026-07-13_2326_part-ii-framework-knob-entry-points-explanation-backward-tracing}
    \item Training-adjustable splits further: into the training loss, into validation, or plain regularizer (details TBD).
    \item Taxonomy = lookup table; the next chapter's method = the procedure that dispatches on it.
    \item UW feature use = the worked class-crossing example: loss term OR up-front feature selection. \vault{2026-07-12_1413_understandability-weighted-feature-use-and-sparsity-axes}
    \item Fidelity vs legibility — the two metric families; every row gets the tag. \vault{2026-07-14_1220_fidelity-vs-legibility-two-metric-families}
    \item Fidelity asks \emph{is the explanation true?} (descriptor-gap measurements); legibility asks \emph{is it simple enough to hold?} \vault{2026-07-14_1218_descriptor-gap-always-a-gap-fidelity-vs-legibility, 2026-07-14_1220_fidelity-vs-legibility-two-metric-families}
    \item Orthogonal by construction: faithful-but-dense and sharp-but-mislabeled both exist.
    \item Legibility caveat: it refers to the \emph{honest model output}, not the post-processed user-facing view; post-processing is explicitly not a concern of this thesis. \vault{2026-07-14_1220_fidelity-vs-legibility-two-metric-families}
    \item Tuning legibility may shrink post-processing AND simplify the internal reasoning itself (careful claim) — possibly even helping performance. \vault{2026-07-13_2108_part-two-motivation-explained-well-tuning-is-reasoning-change}
\end{itemize}
% RELOCATED HERE (ch 2+3 merge, 2026-07-14): the explanation/interpretability
% measurement literature (S4 measuring-why, G1 fidelity metrics) — formerly
% the rw-eval-explanations Related Work section. Cite as a short grounding
% pass where the metric choices are made; no separate survey section.
